\chapter*{Abstract} %l'asterisco dopo chapter serve per visualizzare il capitolo come "non numerato"
La computazione quantistica è un ramo dell'informatica che, con l'avvento dei calcolatori quantistici, è oggi oggetto di intensa ricerca. Un computer quantistico è una macchina che sfrutta i principi e i fenomeni della meccanica quantistica per eseguire calcoli intrattabili per i computer classici. Un campo attraente di questo ramo, ma ancora poco esplorato, è rappresentato dalle strutture dati quantistiche, il cui studio potrebbe portare a significativi miglioramenti nell'efficienza di operazioni fondamentali come la ricerca e l'ordinamento di dati.

L'obiettivo di questa tesi è la costruzione e l'indagine di una struttura dati classica, il Max Heap, nel mondo quantistico e dei suoi potenziali vantaggi rispetto alla controparte classica. Nel primo capitolo sono esposte le basi della computazione quantistica necessarie a comprendere la trattazione successiva. Nel secondo capitolo viene presentato il Quantum Max Heap: se ne discute la costruzione, le operazioni fondamentali e un'analisi dei vantaggi computazionali rispetto al Max Heap classico.